% mcm-table / assets/table-template.tex
% 最小可编译模板 —— 存在的唯一理由是「编译验收需要一个能独立编的最小 .tex」。
% 它**不是样式来源**：表格的规范只在 ../references/table-style.md。
% 自带 booktabs / siunitx；含 \typeout{MCMLAYOUT ...} 供判据取确定的版心宽。
\documentclass[12pt]{article}
\usepackage[margin=1in]{geometry}
\usepackage{booktabs}
\usepackage{siunitx}
\usepackage{amsmath}
\pagestyle{empty}
\typeout{MCMLAYOUT textwidth=\the\textwidth}
\begin{document}

% 来源回显（格式钉死在 ../references/table-style.md §3.2）：放在 \begin{table} 之前，
% 每一格都要有一个 r<行>c<列> 出现在某一行注释的左端。
% mcm-table-src: r1c1 <- pasted:r1c1
% mcm-table-src: r1c2 <- pasted:r1c2
% mcm-table-src: r2c1 <- pasted:r2c1
% mcm-table-src: r2c2 <- pasted:r2c2
% mcm-table-src: r3c1 <- pasted:r3c1
% mcm-table-src: r3c2 <- pasted:r3c2
\begin{table}[htbp]
  \centering
  \caption{An example table typeset with the mcm-table template.}
  \begin{tabular}{l S[table-format=3.2]}
    \toprule
    Item  & {Value} \\
    \midrule
    Alpha & 12.34 \\
    Beta  & {--}  \\
    \bottomrule
  \end{tabular}
\end{table}

\end{document}
